\section{Introducción: por qué este episodio debe leerse como un material de clase trimestral}

Este episodio es la conversación de fin de año del «Informe trimestral de modelos grandes a nivel mundial». No es una entrevista biográfica común, sino una revisión de cara a 2026 de la industria de la IA, de los paradigmas de modelos y de las líneas principales de inversión. Zhang Xiaojun (张小珺) se encarga de llevar las preguntas al plano del negocio, las plataformas y el emprendimiento; Guang Mi (广密), por su parte, condensa un conjunto de observaciones trimestrales en unas cuantas palabras clave: AI War, liderazgo alternado, los modelos fundacionales como plataformas de comercio electrónico generalistas, las dos líneas principales de OpenAI y Google, los tres paradigmas de preentrenamiento/RL/Online Learning y «el modelo es el producto, los datos son el modelo».

Por ello, este material no sigue los subtítulos originales frase por frase, sino que se reorganiza según una lógica didáctica en seis líneas principales. La primera explica por qué «AI Bubble» no basta para explicar esta ronda de inversión: la IA se parece más a una carrera armamentista que nadie puede darse el lujo de perder. La segunda trata las cuentas de ingresos de OpenAI y la disputa entre Google y OpenAI por los puntos de entrada. La tercera aborda el liderazgo alternado de GPT, Claude y Gemini y la analogía entre las empresas de modelos fundacionales y las plataformas. La cuarta trata el tercer paradigma, Online Learning. La quinta abarca los datos, la Agentic Web, los Neo Labs y la Robotics. La sexta trata las narrativas de China y Estados Unidos, los emprendedores de origen chino y las variables de largo plazo para los próximos tres a cinco años.

\begin{importantbox}{Tesis central de este episodio}
El juicio central de Guang Mi no es que «no hay burbuja», sino que «el marco de la burbuja es demasiado estrecho». La inversión en IA está impulsada a la vez por la defensa comercial, la estrategia nacional, la oferta de capacidad de cómputo, la competencia por el talento y las expectativas sobre el siguiente paradigma. Se parece más a una AI War: nadie está seguro de que las cuentas de corto plazo vayan a cuadrar, pero ninguno de los actores clave se atreve a levantarse de la mesa.
\end{importantbox}

\begin{knowledgebox}{Nota sobre la estrategia visual}
Este video no tiene slides, pizarrón, imágenes de artículos ni demostraciones de productos. Conforme al flujo de trabajo para pódcasts, el cuerpo del texto usa la portada solo para identificar la fuente y no vuelve a insertar las tomas fijas de la entrevista; todas las figuras del texto son diagramas conceptuales de elaboración propia que muestran relaciones estructurales, y las explicaciones detalladas se ubican en el texto, las tablas y los recuadros de Lectura de la figura.
\end{knowledgebox}

\subsection{Resumen de la sección}

El valor de EP127 está en condensar los temas de modelos, datos, Agent, capacidad de cómputo y emprendimiento que aparecieron una y otra vez en episodios anteriores en un juicio de fin de año: en 2026 no hay que fijarse en un benchmark aislado, sino en si cambia el paradigma, cómo se redistribuyen los puntos de entrada y quién logra cerrar el ciclo entre datos y producto.
